\documentclass[tikz,border=8pt]{standalone}
\usepackage{ctex}
\usepackage{tikz}
\usetikzlibrary{arrows.meta,positioning}

\begin{document}
\begin{tikzpicture}[
  font=\small,
  >=Latex,
  blk/.style={draw=#1!72!black, fill=#1!8, rounded corners=2mm, very thick, minimum width=30mm, minimum height=10mm, align=center},
  flow/.style={->, very thick, draw=black!75}
]

\node[blk=blue] (c0) at (0,1.6) {\(\mathbf{X}^{can}_t\)\\\(64\times8\)};
\node[blk=teal] (e0) at (0,-1.6) {\(\mathbf{X}^{eth}_t\)\\\(10\times80\)};
\node[blk=purple] (z0) at (4.5,0) {编码并拼接\\\([\mathbf{z}^{can}_t\|\mathbf{z}^{eth}_t]\)};
\node[blk=orange] (seq) at (9.2,0) {\(\mathbf{Z}=[\mathbf{z}_1,\ldots,\mathbf{z}_{10}]\)\\\(10\times d_f\)};
\node[blk=green!70!black] (head) at (14.2,0) {双头输出\\chain + stage};

\draw[flow] (c0.east) -- (z0.west);
\draw[flow] (e0.east) -- (z0.west);
\draw[flow] (z0.east) -- (seq.west);
\draw[flow] (seq.east) -- (head.west);

\node[draw=black!35, fill=blue!7, rounded corners=2mm, align=left, inner sep=4pt] at (7.1,-2.3) {
\textbf{时间构造细节：}\\
每个时间步使用结束对齐 CAN 窗口；\\
当历史帧不足时使用填充策略，帧充足时形成 10 个滑动窗口。};

\end{tikzpicture}
\end{document}
